\documentclass[11pt]{article}
\usepackage[left=0.95in,right=0.95in,top=0.8in,bottom=0.8in]{geometry}
\usepackage{amsmath,amssymb,amsthm,mathtools,booktabs}
\usepackage[T1]{fontenc}
\usepackage{lmodern}
\usepackage[colorlinks=true,linkcolor=blue,urlcolor=blue,citecolor=blue]{hyperref}
\usepackage{microtype}
\newtheorem{theorem}{Theorem}
\newtheorem{proposition}[theorem]{Proposition}
\newtheorem{lemma}[theorem]{Lemma}
\newcommand{\dd}{\,\mathrm{d}}
\newcommand{\HH}[2]{H_{#1}^{(#2)}}
\newcommand{\lcm}{\operatorname{lcm}}
\newcommand{\Z}{\mathbb Z}
\newcommand{\Q}{\mathbb Q}
\newcommand{\R}{\mathbb R}
\title{Odd-zeta linear forms and their arithmetic cost}
\author{A derivational note accompanying the Lean development}
\date{27 October 2026}
\begin{document}
\maketitle
\begin{abstract}
A positive moment determinant isolates any fixed odd zeta value without a
weight projector. We derive its nonvanishing, construct explicit integer
coefficients, and account for denominator clearing. For a fixed moment shift,
the raw determinant has logarithmic decay rate $\log 4$ on the square of its
size. An elementary binomial divisibility bound shows that the uniform
least-common-multiple normalization instead grows at every integer weight
$s>1$. Exact common-factor cancellation is therefore essential for this
family. The accompanying Lean file verifies the integer construction and
its conditional irrationality criterion. No unconditional irrationality
theorem for all odd zeta values is established.
\end{abstract}

\section{The quantity an irrationality argument must make small}
Write a rational linear form as $L=A\alpha+B$, and choose an integer
$D>0$ such that $a=DA$ and $b=DB$ are integers. If $(a,b)\ne(0,0)$, set
$g=\gcd(|a|,|b|)$ and
\begin{equation}
 \ell=\frac{D}{g}L=\frac ag\alpha+\frac bg .
 \label{eq:primitive}
\end{equation}
Under the hypothesis $\alpha=p/q\in\Q$, every nonzero $\ell$ satisfies
$|q\ell|\ge1$. Thus arbitrarily late nonzero forms tending to zero imply
irrationality. The exact arithmetic budget is
\begin{equation}
 -\log|\ell|=-\log|L|-\log D+\log g .
 \label{eq:budget}
\end{equation}
An analytic estimate on $L$ alone does not control this budget.
The common divisor in \eqref{eq:primitive} must be proved, rather than
inferred from a rational or geometric change of coordinates.

\section{A pure-target construction at every integer weight}
Fix integers $s>1$, $a\ge0$, and $N\ge1$. Define
\begin{equation}
 \HH{k}{s}=\sum_{j=1}^k j^{-s},\qquad
 \mu_k=\zeta(s)-\HH{k}{s},\qquad
 \Delta_{N,a}^{(s)}=\det[\mu_{a+i+j}]_{0\le i,j<N}.
 \label{eq:moment}
\end{equation}
The case requested in the chat is $s=2r+1$ with $r\ge1$.

\begin{theorem}
For these parameters there are explicit $A_{N,a}^{(s)},B_{N,a}^{(s)}\in\Q$
such that
\begin{equation}
 0<\Delta_{N,a}^{(s)}
   =A_{N,a}^{(s)}\zeta(s)+B_{N,a}^{(s)}.
 \label{eq:affine}
\end{equation}
In particular, this construction isolates one target and is nonzero.
\end{theorem}
\begin{proof}
Substitute $x=e^{-t}$ in the elementary gamma integral to obtain
\[
 \frac1{(j+1)^s}
 =\frac1{(s-1)!}\int_0^1 x^j(-\log x)^{s-1}\dd x .
\]
Summing nonnegative integrands gives
\begin{equation}
 \mu_k=\int_0^1 x^k w_s(x)\dd x,\qquad
 w_s(x)=\frac{(-\log x)^{s-1}}{(s-1)!(1-x)} .
 \label{eq:weight}
\end{equation}
This measure is finite: its total mass is $\zeta(s)<\infty$.
Expanding the determinant and symmetrizing its product integrals gives
\begin{equation}
 \Delta_{N,a}^{(s)}
 =\frac1{N!}\int_{(0,1)^N}
  \prod_{i<j}(x_i-x_j)^2
  \prod_{i=1}^N x_i^a w_s(x_i)\dd x_i>0 .
 \label{eq:gram}
\end{equation}
The integrand is bounded above by the product of the finite measures.
It is strictly positive when the variables lie in pairwise disjoint
closed interior intervals, proving strict positivity of the integral.

For the affine identity, put
\[
 P(x)=\det[x-\HH{a+i+j}{s}]_{0\le i,j<N}.
\]
Subtract row zero from every other row. Only row zero then depends on $x$,
so determinant multilinearity gives
$P(x)=(P(1)-P(0))x+P(0)$. Take $A=P(1)-P(0)$ and $B=P(0)$.
\end{proof}

The coefficient of the target has a useful independent expression:
\begin{equation}
 A_{N,a}^{(s)}
 =\det\left[
 \frac1{(a+i+j+1)^s}-\frac1{(a+i+j+2)^s}
 \right]_{0\le i,j<N-1}>0 .
 \label{eq:coefficient}
\end{equation}
Here an empty determinant is $1$. For $N=1$ the coefficient is $1$ directly.
For $N\ge2$, replace the monomial basis by
\[
 1,\quad(1-x),\quad x(1-x),\quad\ldots,\quad x^{N-2}(1-x).
\]
This integer basis change has determinant
$\pm1$. Only the first polynomial is nonzero at $1$, so the coefficient of
$\zeta(s)$ is the lower principal minor. Its entries are
\[
 \mu_k-2\mu_{k+1}+\mu_{k+2}=(k+1)^{-s}-(k+2)^{-s}.
\]
The same positive Gram argument proves the inequality.

\section{Explicit denominator clearing, including its cost}
Let $M=a+2N-2$, $d_M=\lcm(1,\ldots,M)$, with $d_0=1$, and put
$Q=d_M^s$. For every $m\le M$,
\begin{equation}
 Q\HH{m}{s}
 =\sum_{j=1}^m(d_M/j)^s\in\Z .
 \label{eq:hclear}
\end{equation}
Define the integer matrix and two integers
\begin{align}
 Z_{ij}&=-\sum_{j'=1}^{a+i+j}(d_M/j')^s,\\
 U&=\det[Q+Z_{ij}]-\det[Z_{ij}],\qquad
 V=\det[Z_{ij}].
 \label{eq:integers}
\end{align}
Applying the affine determinant identity after scaling gives the exact relation
\begin{equation}
 \boxed{Q^N\Delta_{N,a}^{(s)}=U\zeta(s)+V,\qquad U,V\in\Z.}
 \label{eq:integerform}
\end{equation}
Consequently the primitive form is
$(Q^N/g)\Delta_{N,a}^{(s)}$, where $g=\gcd(|U|,|V|)$.
These definitions, including \eqref{eq:hclear} and \eqref{eq:integerform},
are proved in the accompanying Lean module for every natural $s$.
The positivity statement uses $s>1$ and is proved above in ordinary mathematics.

For $M\ge1$, its exact prime cost is $v_p(Q^N)=sN\lfloor\log_p M\rfloor$.
The scale $Q^N$ is safe, but it is expensive. To quantify this without the
prime number theorem, observe that
\begin{equation}
 \binom{M}{k}\mid d_M\quad(0\le k\le M),\qquad
 2^M\le(M+1)d_M .
 \label{eq:lcm}
\end{equation}
The case $M=0$ is immediate. For $M\ge1$, Legendre's formula gives, for every prime $p$,
\[
 v_p\binom{M}{k}
 =\sum_{\nu\ge1}
 \left(\left\lfloor\frac M{p^\nu}\right\rfloor
 -\left\lfloor\frac k{p^\nu}\right\rfloor
 -\left\lfloor\frac{M-k}{p^\nu}\right\rfloor\right)
 \le\lfloor\log_p M\rfloor=v_p(d_M).
\]
Each summand is $0$ or $1$ and is zero for $p^\nu>M$.
Summing the $M+1$ binomial coefficients proves the second inequality.
It follows that
\begin{equation}
 \log(Q^N)\ge sN\bigl(M\log2-\log(M+1)\bigr).
 \label{eq:cost}
\end{equation}

\section{The exact analytic scale and a normalization obstruction}
\begin{theorem}
For every fixed integer $s>1$ and fixed $a\ge0$,
\begin{equation}
 \lim_{N\to\infty}\frac{\log\Delta_{N,a}^{(s)}}{N^2}=-\log4 .
 \label{eq:rate}
\end{equation}
Moreover the uniform integer normalization in \eqref{eq:integerform} satisfies
\begin{equation}
 \liminf_{N\to\infty}
 \frac{\log(Q^N\Delta_{N,a}^{(s)})}{N^2}
 \ge(s-1)\log4>0.
 \label{eq:obstruction}
\end{equation}
Thus this normalization cannot be the small-form sequence.
\end{theorem}
\begin{proof}
Let $\rho(x)=x^a w_s(x)$. Use $p_0=1$ and the monic shifted Chebyshev polynomials
\[
 p_j(x)=2^{1-2j}T_j(2x-1)\qquad(j\ge1).
\]
The monic basis change preserves the Gram determinant. Hadamard's inequality,
$|T_j|\le1$ on $[-1,1]$, and $\int\rho=\mu_a$ yield
\begin{equation}
 0<\Delta_{N,a}^{(s)}
 \le\prod_{j=0}^{N-1}\int_0^1p_j(x)^2\rho(x)\dd x
 \le\mu_a^N\,2^{-2(N-1)^2}.
 \label{eq:upper}
\end{equation}
This proves the required upper logarithmic limit.

For the lower bound, fix $0<\varepsilon<1/2$, let
$\lambda=1-2\varepsilon$, and let
$c_\varepsilon=\min_{[\varepsilon,1-\varepsilon]}\rho>0$.
Restricting \eqref{eq:gram} to this interval and making an affine substitution
gives
\begin{equation}
 \Delta_{N,a}^{(s)}
 \ge c_\varepsilon^N\lambda^{N^2} J_N,\qquad
 J_N=\det[(i+j+1)^{-1}]_{i,j<N}.
 \label{eq:lower}
\end{equation}
The Cauchy determinant formula, obtained by clearing denominators and
factoring its alternating numerator, gives
\[
 J_N=\prod_{j=0}^{N-1}
 \frac{(j!)^4}{(2j)!(2j+1)!}
 =\prod_{j=0}^{N-1}
 \frac1{(2j+1)\binom{2j}{j}^{\,2}} .
\]
Since $4^j/(2j+1)\le\binom{2j}{j}\le4^j$, the logarithm of this product is
$-N(N-1)\log4+O(N\log(N+1))$. Hence
$\liminf N^{-2}\log\Delta\ge\log\lambda-\log4$.
Let $\varepsilon$ decrease to zero to prove \eqref{eq:rate}.

Finally, $M=a+2N-2$ and \eqref{eq:cost} imply
$\liminf N^{-2}\log(Q^N)\ge s\log4$.
Adding \eqref{eq:rate} proves \eqref{eq:obstruction}.
\end{proof}

This obstruction concerns the uniform scale $Q^N$. It does not bound the
primitive common divisor $g$. For this fixed-shift family to yield
primitive forms tending to zero, the exact requirement is
\begin{equation}
 \log g-\log(Q^N)-\log\Delta_{N,a}^{(s)}\longrightarrow+\infty.
 \label{eq:remaining}
\end{equation}
In particular, \eqref{eq:obstruction} forces a saving of at least
$(s-1)\log4$ on the $N^2$ scale. A rational change of polynomial basis in
the same space multiplies both coefficients by $(\det T)^2$; after
primitive reduction it gives the same integer pair up to sign.
Improving the budget therefore requires additional arithmetic identities
or a different polynomial space or measure. A shift $a=a_N$ requires a
new asymptotic calculation; theorem \eqref{eq:rate} assumes $a$ fixed.

\section{Checks, source connections, and formalization boundary}
The exact rational audit already in the repository clears reduced
denominators and removes the entire gcd. At $s=5$, $N=20$, $a=0$, it encloses
the primitive form strictly above $1$ in absolute value. The displayed rates
$-\log|\Delta|/N^2\approx1.517981$ and
$-\log|\ell|/N^2\approx-4.772732$ are finite-index diagnostics.
They are not used in the proofs above and do not exclude eventual
arithmetic improvements.

An additional exact search replaces the monomials by sparse exponents.
It checks $700$ small exponent sets and $60$ larger consecutive, even,
quadratic, gapped, and two-cluster examples across $s=3,5,7,9$.
Every tested example of the five families at sizes $4,6,8$ has primitive
absolute value greater than $1$. These are finite checks, not an
impossibility theorem for sparse spaces. The source chat's earlier weight
projection and its separate sign obstruction are recorded in the repository audit.

Recent geometric constructions are useful specifically for their
denominator symmetries. Brown--Zudilin \cite{BZ} exploit a group of order
$5040$ for cellular integrals, obtaining an effective approximation
exponent $0.86$ at weight five. Their result does not prove irrationality.
The $2$-adic construction of Lai--Sprang--Zudilin \cite{LSZ} concerns
$\zeta_2(5)$; its smallness occurs in a different absolute value.
Neither result supplies the missing real estimate \eqref{eq:remaining}
for the present determinant family.

\medskip
\noindent\textbf{What the Lean file verifies.}
\texttt{OddHankelArithmetic.lean} defines the harmonic sums, the lcm,
the integer matrix, and $U,V$. It proves \eqref{eq:hclear},
\eqref{eq:integerform}, the lcm lower bound, the finite cost \eqref{eq:cost},
and its exact prime valuations. It also verifies division by any proved
common divisor and the resulting conditional irrationality criterion.
The moment integrals, positivity, and asymptotics
\eqref{eq:rate}--\eqref{eq:obstruction} are paper proofs here, not
Lean-verified results. No unconditional all-odd irrationality conclusion follows.

\begin{thebibliography}{9}
\bibitem{BZ}
F. Brown and W. Zudilin,
\emph{On cellular rational approximations to $\zeta(5)$},
arXiv:2210.03391v3, revised 29 January 2026.
\url{https://arxiv.org/abs/2210.03391v3}.
\bibitem{LSZ}
L. Lai, J. Sprang and W. Zudilin,
\emph{A note on the irrationality of $\zeta_2(5)$},
International Mathematics Research Notices 2026, rnag180.
\url{https://arxiv.org/abs/2505.05005v2}.
\end{thebibliography}
\end{document}
